\subsection{Complexes of homomorphisms and homotopies}

PROPOSITION 3. --- Let C, \(\tilde{C}\) , \(C'\) , \(\tilde{C}'\) be four A-complexes, \(u: \tilde{C} \to C\) , \(v: \tilde{C} \to C\) , \(u': C' \to \tilde{C}'\) and \(v': C' \to \tilde{C}'\) four morphisms of complexes.

\begin{enumerate}
\def\labelenumi{\alph{enumi})}
\item
  If u and \(u'\) are homotopic to v and \(v'\) respectively, then the two morphisms \(\operatorname{Homgr}_{\mathbf{A}}(u, u')\) and \(\operatorname{Homgr}_{\mathbf{A}}(v, v')\) of \(\operatorname{Homgr}_{\mathbf{A}}(\mathbf{C}, \mathbf{C}')\) to \(\operatorname{Homgr}_{\mathbf{A}}(\tilde{\mathbf{C}}, \tilde{\mathbf{C}}')\) are homotopic.
\item
  If u and u' are homotopisms, Homgr\_A(u, u') is a homotopism.
\item
  If C or C' is homotopic to zero, Homgr\_A(C, C') is homotopic to zero.
\end{enumerate}

Let us denote by the same letter \(d\) the differentials of the complexes C, C \(_{1}\) , C \(^{\prime}\) , C \(_{1}^{\prime}\) , and by D the differentials of Homgr \(_{A}\) (C, C') and Homgr \(_{A}\) (C \(_{1}\) , C \(_{1}^{\prime}\) ). If \(u\) (resp. \(u^{\prime}\) ) is homotopic to \(v\) (resp. \(v^{\prime}\) ), there exists a graded homomorphism of degree 1,

\[
w: \mathrm{C} _ {1} \rightarrow \mathrm{C} \quad (\text { resp. } w ^ {\prime}: \mathrm{C} ^ {\prime} \rightarrow \mathrm{C} _ {1} ^ {\prime})
\]

such that

\[
u - v = d w + w d \quad (\text { resp. } u ^ {\prime} - v ^ {\prime} = d w ^ {\prime} + w ^ {\prime}\tag{2}
\]

Let W : Homgr \(_{A}\) (C, C') → Homgr \(_{A}\) (C \(_{1}\) , C \(_{1}\) ) be the graded homomorphism of degree 1 such that, for \(f \in\) Homgr \(_{A}\) (C, C') \(_{n}\) , \(n \in\) Z, we have

\[
\mathrm{W} (f) = w ^ {\prime} f u + (- 1) ^ {n} v ^ {\prime} f w.\tag{3}
\]

We then have

\[
\begin{array}{r l} (\mathrm{DW} + \mathrm{WD}) (f) & = \mathrm{D} [ w ^ {\prime} f u + (- 1) ^ {n} v ^ {\prime} f w ] + \mathrm{W} [ d f - (- 1) ^ {n} f d ] \\ & = d w ^ {\prime} f u - (- 1) ^ {n + 1} w ^ {\prime} f u d + (- 1) ^ {n} d v ^ {\prime} f w + v ^ {\prime} f w d \\ & \quad + w ^ {\prime} d f u + (- 1) ^ {n + 1} v ^ {\prime} d f w - (- 1) ^ {n} w ^ {\prime} f d u + v ^ {\prime} f d w \\ & = (d w ^ {\prime} + w ^ {\prime} d) f u + v ^ {\prime} f (w d + d w) \\ & = (u ^ {\prime} - v ^ {\prime}) f u + v ^ {\prime} f (u - v) = u ^ {\prime} f u - v ^ {\prime} f v; \end{array}
\]

which is written DW + WD = Homgr\_A (u, u')-Homgr\_A (v, v'), whence a).

Let us prove b). If u and \(u'\) are homotopisms, let \(\alpha: C \to \tilde{C}\) and \(\alpha': \tilde{C}' \to C'\) be morphisms of complexes such that \(u \circ \alpha\) , \(\alpha \circ u\) , \(u' \circ \alpha'\) , \(\alpha' \circ u'\) are homotopic respectively to \(Id_{C}\) , \(Id_{\tilde{C}}\) , \(Id_{\tilde{C}'}\) , \(Id_{C'}\) . Then Homgr \((u, u') \circ \text{Homgr } (\alpha, \alpha')\) , which is equal to Homgr \((\alpha \circ u, u' \circ \alpha')\) , is homotopic by a) to

\[
\operatorname{Homgr} _ {A} \left(\operatorname{Id} _ {C}, \operatorname{Id} _ {C ^ {\prime}}\right) = \operatorname{Id} _ {\operatorname{Homgr} (C, C ^ {\prime})}:
\]

likewise Homgr \((\alpha, \alpha') \circ\) Homgr \((u, u')\) is homotopic to \(\operatorname{Id}_{\operatorname{Homgr}(\tilde{\mathbb{C}}, \tilde{\mathbb{C}})}\) , whence \(b\) ).

Finally c) follows from b) (applied to the case where \(\tilde{C}\) or \(\tilde{C}'\) is zero).

COROLLARY 1. --- If C is split and H(C) is projective (resp. if C' is split and H \(_{n}\) (C') injective for each n), then the canonical homomorphism

\[
\lambda (\mathrm{C}, \mathrm{C} ^ {\prime}): \mathrm{H} (\operatorname{Homgr} _ {\mathrm{A}} (\mathrm{C}, \mathrm{C} ^ {\prime})) \rightarrow \operatorname{Homgr} _ {\mathrm{A}} (\mathrm{H} (\mathrm{C}), \mathrm{H} (\mathrm{C} ^ {\prime}))
\]

is bijective.

Suppose for example that C' is split and H(C') is injective for each n, the case where C is split and H(C) is projective being proved analogously. By X, p. 35, def. 6, there exists a homotopism \(u': C' \to H(C')\) ; by prop. 3, Homgr\_A (1, \(u'\) ) is a homotopism of Homgr\_A (C, C') on Homgr\_A (C, H(C')); since

\[
\operatorname{Homgr} _ {\mathbf {A}} \left(1, \mathrm{H} \left(u ^ {\prime}\right)\right) \circ \lambda (\mathrm{C}, \mathrm{C} ^ {\prime}) = \lambda (\mathrm{C}, \mathrm{H} \left(\mathrm{C} ^ {\prime}\right)) \circ \mathrm{H} \left(\operatorname{Homgr} _ {\mathbf {A}} \left(1, u ^ {\prime}\right)\right)
\]

and that \(\operatorname{Homgr}_{\mathbf{A}}(1, \mathrm{H}(u'))\) and \(\operatorname{H}(\operatorname{Homgr}_{\mathbf{A}}(1, u'))\) are bijective, it suffices for us to prove that \(\lambda(\mathbf{C},\mathbf{H}(\mathbf{C}^{\prime}))\) is bijective, which brings us back to the case where \(\mathbf{C}'\) is injective and has zero differential.

Then the canonical exact sequences (X, p. 28)

\[
0 \rightarrow \mathrm{B} (\mathrm{C}) \xrightarrow {i} \mathrm{C} \xrightarrow {p} \mathrm{C} / \mathrm{B} (\mathrm{C}) \rightarrow 0\tag{III}
\]

\[
0 \rightarrow \mathrm{H} (\mathrm{C}) \stackrel {{j}} {{\rightarrow}} \mathrm{C} / \mathrm{B} (\mathrm{C}) \stackrel {{\delta}} {{\rightarrow}} \mathrm{B} (\mathrm{C}) \rightarrow 0\tag{IV}
\]

give exact sequences (X, p. 83, prop. 2, a))

\[
0 \rightarrow \operatorname{Homgr} _ {\mathrm{A}} (\mathrm{C} / \mathrm{B} (\mathrm{C}), \mathrm{C} ^ {\prime}) \xrightarrow {\mathrm{P}} \operatorname{Homgr} _ {\mathrm{A}} (\mathrm{C}, \mathrm{C} ^ {\prime}) \xrightarrow {\mathrm{I}} \operatorname{Homgr} _ {\mathrm{A}} (\mathrm{B} (\mathrm{C}), \mathrm{C} ^ {\prime}) \rightarrow 0
\]

\[
0 \rightarrow \operatorname{Homgr} _ {\mathrm{A}} (\mathrm{B} (\mathrm{C}), \mathrm{C} ^ {\prime}) \stackrel {{\Delta}} {{\rightarrow}} \operatorname{Homgr} _ {\mathrm{A}} (\mathrm{C} / \mathrm{B} (\mathrm{C}), \mathrm{C} ^ {\prime}) \stackrel {{\mathrm{J}}} {{\rightarrow}} \operatorname{Homgr} _ {\mathrm{A}} (\mathrm{H} (\mathrm{C}), \mathrm{C} ^ {\prime}) \rightarrow 0.
\]

Since \(d_{\mathbf{C}} = i \circ \delta \circ p\) the differential D of Homgr(C, C') is given by \(\mathrm{D}_{n} = (-1)^{n+1} \mathrm{P}_{n} \circ \Delta_{n} \circ \mathrm{I}_{n}\) ; we then have

\[
\mathrm{Z} \left(\operatorname{Homgr} _ {\mathrm{A}} \left(\mathrm{C}, \mathrm{C} ^ {\prime}\right)\right) = \operatorname{Ker} (\mathrm{P} \circ \Delta \circ \mathrm{I}) = \operatorname{Ker} \mathrm{I} = \operatorname{Im} \mathrm{P},
\]

\[
\mathrm{B} \left(\operatorname{Homgr} _ {\mathrm{A}} \left(\mathrm{C}, \mathrm{C} ^ {\prime}\right)\right) = \operatorname{Im} (\mathrm{P} \circ \Delta \circ \mathrm{I}) = \mathrm{P} (\operatorname{Im} \Delta) = \mathrm{P} (\text { Ker   J });
\]

whence an isomorphism \(\varphi : H(\text{Homgr}_A(C, C')) \to \text{Homgr}_A(H(C), C')\) such that, if \(a \in \text{Homgr}_A(C/B(C), C')\) , then the image under \(\varphi\) of the class of \(P(a)\) is \(J(a)\) ; one immediately verifies that \(\varphi\) identifies with the canonical homomorphism \(\lambda\) .

COROLLARY 2. --- Suppose B(C) and H(C) are projective (resp. B \(_{n}\) (C') and H \(_{n}\) (C') injective for each n). Then λ(C, C') is bijective.

Indeed C (resp. C′) is then split by X, p. 35, example 4, and we apply Cor. 1.

COROLLARY 3. --- Let M be a projective (resp. injective) A-module. For every complex C of A-modules and every integer n, the canonical homomorphism

\[
\mathrm{H} ^ {n} (\operatorname{Homgr} _ {\mathbf {A}} (\mathrm{M}, \mathrm{C})) \rightarrow \operatorname{Hom} _ {\mathbf {A}} (\mathrm{M}, \mathrm{H} ^ {n} (\mathrm{C}))
\]

\[
\left(\text { resp. } \mathrm{H} ^ {n} (\operatorname{Homgr} _ {\mathbf {A}} (\mathbf {C}, \mathbf {M})) \rightarrow \operatorname{Hom} _ {\mathbf {A}} \left(\mathrm{H} _ {n} (\mathbf {C}), \mathbf {M}\right)\right) \text { is   bijective. }
\]

Lemma 1. --- a) If C or C′ is bounded on the right, if C is projective and if H(C′) = 0, then H(Homgr\_A(C, C')) = 0.

\begin{enumerate}
\def\labelenumi{\alph{enumi})}
\setcounter{enumi}{1}
\tightlist
\item
  If C or C' is bounded on the left, if C' is injective and if H(C) = 0, then H(Homgr\_A(C, C')) = 0.
\end{enumerate}

Let \(f \in Z_n(\text{Homgr}_A(C, C'))\) ; \(f\) is therefore a morphism of complexes of \(C\) to \(C'(n)\) ; in the case \(a\) (resp. \(b\) ), \(f_m\) is zero for \(m\) small enough (resp. large enough). By \(X\) , p. 47, prop. 1, \(f\) is then homotopic to zero, hence belongs to \(B_n(\text{Homgr}_A(C, C'))\) whence the conclusion.

PROPOSITION 4. --- Let \(u: \mathbb{C}' \to \mathbb{C}\) be a homologism of complexes, \(\mathbb{P}\) a projective complex, \(\mathbb{E}\) an injective complex.

\begin{enumerate}
\def\labelenumi{\alph{enumi})}
\tightlist
\item
  If P is bounded above, or if C and C' are bounded above, then
\end{enumerate}

\[
\operatorname{Homgr} _ {\mathbf {A}} (1, u): \operatorname{Homgr} _ {\mathbf {A}} (\mathrm{P}, \mathrm{C} ^ {\prime}) \rightarrow \operatorname{Homgr} _ {\mathbf {A}} (\mathrm{P}, \mathrm{C})
\]

is a homologism.

\begin{enumerate}
\def\labelenumi{\alph{enumi})}
\setcounter{enumi}{1}
\tightlist
\item
  If E is bounded below, or if C and C' are bounded below, then
\end{enumerate}

\[
\operatorname{Homgr} _ {\mathbf {A}} (u, 1): \operatorname{Homgr} _ {\mathbf {A}} (\mathrm{C}, \mathrm{E}) \rightarrow \operatorname{Homgr} _ {\mathbf {A}} (\mathrm{C} ^ {\prime}, \mathrm{E})
\]

is a homologism.

Suppose first that u is injective, and set \(C'' = \text{Coker } u\) . Since u is a homologism, \(C''\) has zero homology. On the other hand, we have exact sequences (prop. 2)

\[
0 \rightarrow \operatorname{Homgr} _ {\mathrm{A}} (\mathrm{P}, \mathrm{C} ^ {\prime}) \xrightarrow {\operatorname{Homgr} (1 , u)} \operatorname{Homgr} _ {\mathrm{A}} (\mathrm{P}, \mathrm{C}) \rightarrow \operatorname{Homgr} _ {\mathrm{A}} (\mathrm{P}, \mathrm{C} ^ {\prime \prime}) \rightarrow 0
\]

\[
0 \rightarrow \operatorname{Homgr} _ {\mathbf {A}} \left(\mathrm{C} ^ {\prime \prime}, \mathrm{E}\right)\rightarrow \operatorname{Homgr} _ {\mathbf {A}} (\mathrm{C}, \mathrm{E}) \xrightarrow {\operatorname{Homgr} (u , 1)} \operatorname{Homgr} _ {\mathbf {A}} \left(\mathrm{C} ^ {\prime}, \mathrm{E}\right)\rightarrow 0
\]

by lemma 1, Homgr \(_{A}\) (P, C'') has zero homology in the case \(a\) ), Homgr \(_{A}\) (C'', E) has zero homology in the case \(b\) ), whence the conclusion.

In the general case, there exists (X, p. 38, cor. to prop. 7) a complex \(\tilde{\mathbf{C}}'\) , which is bounded on the right (resp. on the left) when C and C' are, an injective morphism \(\tilde{u}: \mathbf{C}' \to \tilde{\mathbf{C}}'\) and a homotopism \(\beta: \tilde{\mathbf{C}}' \to \mathbf{C}\) such that \(u = \beta \circ \tilde{u}\) . Then \(\tilde{u}\) is a homologism (X, p. 34, cor. to prop. 5); by the preceding, Homgr \(_{\mathbf{A}}\) ( \(1_{\mathbf{P}}, \tilde{u}\) ) (resp. Homgr \(_{\mathbf{A}}\) ( \(\tilde{u}, 1_{\mathbf{E}}\) )) is a homologism in the case \(a\) ) (resp. \(b\) ). On the other hand Homgr \(_{\mathbf{A}}\) ( \(1_{\mathbf{P}}, \beta\) ) (resp. Homgr \(_{\mathbf{A}}\) ( \(\beta, 1_{\mathbf{E}}\) )) is a homotopism (prop. 3): then Homgr \(_{\mathbf{A}}\) ( \(1_{\mathbf{P}}, u\) ) (resp. Homgr \(_{\mathbf{A}}\) ( \(u, 1_{\mathbf{E}}\) )) is composed of two homologisms, hence is a homologism.

